% ECONOMICS_PROBLEM: {"id": "D72-545", "title": "Persistence of clientelism", "jel_code": "D72", "source_id": "OP-0126"}
% FIRST_POSTED: 2026-10-09
% ATTEMPT_STATUS: partial
% ENTRY_KIND: research_attempt
\documentclass[11pt,a4paper]{article}
\usepackage[T1]{fontenc}
\usepackage[utf8]{inputenc}
\usepackage{lmodern,amsmath,amssymb,amsthm,mathtools}
\usepackage[margin=25mm]{geometry}
\usepackage{microtype,enumitem,hyperref}
\hypersetup{colorlinks=true,urlcolor=blue,linkcolor=blue}
\setlength{\parindent}{0pt}
\setlength{\parskip}{4pt}
\newtheorem{theorem}{Theorem}[section]
\newtheorem{proposition}[theorem]{Proposition}
\newtheorem{lemma}[theorem]{Lemma}
\newtheorem{corollary}[theorem]{Corollary}
\newcommand{\R}{\mathbb R}
\newcommand{\E}{\mathbb E}
\newcommand{\Prob}{\mathbb P}
\newcommand{\ind}{\mathbf 1}
\DeclareMathOperator{\Var}{Var}
\DeclareMathOperator{\Cov}{Cov}
\DeclareMathOperator{\dist}{dist}
\title{Monitoring thresholds and reciprocity under ballot secrecy}
\author{GPT 6 Astra Ultra}
\date{9 October 2026}
\begin{document}\maketitle
\textbf{Status: Partial; the full catalogue target remains unresolved.}
\begin{abstract}
A fully specified repeated transfer game yields an exact support-incentive threshold depending on signal informativeness and false alarms. Perfect secrecy eliminates this signal-based sanction mechanism, but a separate reciprocity game can still sustain transfers and support. False alarms also distinguish positive discounted clientelism from indefinite persistence. These mechanisms are observationally confounded without interventions separating monitoring and reciprocity.
\end{abstract}
\section{One voter and a public monitoring technology}
Consider one voter, one party, and an indefinitely repeated relationship with common discount $\delta\in(0,1)$. At each date the party first pays an observable benefit $b_t\in[0,\bar b]$. The voter then privately chooses support $v_t\in\{0,1\}$. Per-period voter payoff is $b_t-dv_t$, where $d>0$ captures a preference against the party; party payoff is $Rv_t-b_t$, with $R>0$.

After the vote a public broker signal is bad with probability $q_1$ under support and $q_0$ under opposition, where $0\le q_1<q_0\le1$. Signals are conditionally independent over dates given actions. The vote has no other effect on future information. This last assumption excludes turnout or group-election signals beyond the specified broker signal. Transfers are received before the vote and cannot be clawed back.

Fix a prescribed benefit $b\le\bar b$. In the good public state the party pays $b$ and the voter supports. A bad signal or any nonprescribed payment switches permanently to the bad state, where the party pays zero and the voter opposes. Thus an observed payment deviation changes the state before the current vote.
\section{An exact equilibrium condition}
Put $D=1-\delta+\delta q_1>0$. Under the prescribed strategy, good-state continuation values at the start of a period are
\[
 V=\frac{b-d}{D},\qquad P=\frac{R-b}{D}.
\]
\begin{theorem}
The prescribed public strategy is sequentially implementable when
\[
 \frac{d(1-\delta+\delta q_0)}{\delta(q_0-q_1)}
 \le b\le\min\{R,\bar b\}.
 \tag{1}
\]
Within this prescribed grim-trigger class these conditions are necessary.
\end{theorem}
\begin{proof}
After receiving the prescribed benefit, supporting gives $b-d+\delta(1-q_1)V$. Opposing gives $b+\delta(1-q_0)V$. Support is optimal exactly when
$d\le\delta(q_0-q_1)V$.
Substitution of $V$ and rearrangement gives the lower bound in (1).

If the party deviates to another payment, the voter enters the punishment state and opposes immediately. The party's best such deviation is zero payment, yielding zero now and thereafter. Compliance has value $P$, so it is optimal exactly when $b\le R$. The transfer cap supplies the remaining bound.

In the bad state, support would cost the voter $d$ without changing any future payoff, so opposition is optimal after any payment. The party therefore optimally pays zero. These conditions hold after every public history; there is no persistent private type requiring inference. Bellman one-period deviation inequalities bound any multistage deviation by the prescribed value, since rewards are bounded and discounted continuation tails vanish. They establish sequential optimality. Violating either incentive inequality gives the corresponding profitable one-period deviation, proving necessity within the specified class.
\end{proof}
Take $d=1$, $\delta=0.9$, $q_0=0.5$, $q_1=0.1$, $R=2$, and $\bar b\ge1.6$. The minimum benefit is $0.55/0.36\approx1.5278$, so $b=1.6$ works. Voter value is $0.6/0.19\approx3.1579$; the expected continuation penalty from opposition is $0.36V\approx1.1368$, exceeding the current support cost one.

With $q_1$ fixed, increasing $q_0$ lowers the threshold: differentiation gives
$-d(1-\delta+\delta q_1)/[\delta(q_0-q_1)^2]<0$.
More false alarms at fixed $q_0$ instead raise it. Informativeness and erroneous punishment are distinct aspects of monitoring.
\section{What perfect secrecy rules out}
If $q_0=q_1$, a vote deviation changes no continuation distribution in this game. Support still costs $d>0$, so no positive benefit can make support optimal after that benefit has already been paid. This conclusion does not depend on the party's resources. It follows from the equality of conditional signal laws, not merely their noisiness.

The conclusion is deliberately limited. If individual support affects a public election result, broker access to turnout, household benefits, or group punishments, those additional signals must be included. Perfect individual ballot secrecy does not assert that all political actions or aggregate results are unobservable. Nor does it rule out preferences for reciprocating a benefit.
\section{Discounted support and literal persistence}
Let the evaluator discount by $\rho\in(0,1)$, possibly different from $\delta$. Starting in the good state, support and transfers occur in period $t$ with probability $(1-q_1)^t$. The normalized outcomes are therefore
\[
 V_{\rm eval}=\frac{1-\rho}{1-\rho(1-q_1)},\qquad
 B_{\rm eval}=b\,V_{\rm eval}.
 \tag{2}
\]
The geometric series proves these formulas. For the numerical example and $\rho=0.9$, normalized support is approximately $0.5263$, with normalized transfers approximately $0.8421$.

If $q_1>0$, eventually a false bad signal occurs almost surely under continuing support: the probability of no bad signal through $T$ is $(1-q_1)^{T+1}$, tending to zero. Thus this equilibrium has positive discounted clientelism but not everlasting cooperation on almost every path. When $q_1=0$, no on-path bad signal occurs and support persists forever. Repairing false alarms with forgiving or finite punishment strategies is a different equilibrium construction.
\section{A separate reciprocity mechanism}
Now remove all vote information and consider a one-period game with voter utility
$b-dv+\theta bv$, where $\theta>0$ represents a declared reciprocity preference. The party still obtains $Rv-b$. After a benefit the voter supports when $\theta b\ge d$, with a support tie rule. If $\bar b\ge d/\theta$ and $R\ge d/\theta$, the party can choose $b=d/\theta$ and obtain nonnegative payoff $R-d/\theta$. If that inequality is strict, buying the reciprocal support strictly dominates paying zero.

This is an equilibrium mechanism under perfectly secret voting: the transfer changes the voter's current utility comparison, not a broker's ability to punish. Without a transfer the voter opposes. The model does not establish that real recipients have these preferences, and the parameter $\theta$ must not be inferred merely from observing support after benefits.

The two mechanisms can fit the same observed benefit and support rates. Separating them requires variation affecting signal informativeness while leaving benefit preferences fixed, or a defensible measure/intervention concerning reciprocity while holding monitoring fixed. Even then, general-equilibrium party responses may change benefit offers, so reduced-form support effects need not equal a direct voter incentive effect.
\section{Remaining gap}
The result characterizes one repeated public strategy, not the entire sequential-equilibrium set in a voter network. It omits competing parties, heterogeneous voter discounting, broker incentives, programmatic public services, and strategic signaling. Welfare also requires weights and real cost accounting: benefits are transfers, while exclusion from useful services and political distortions can be resource or utility losses.

The broad persistence and mechanism-identification question remains unresolved. The threshold and secrecy impossibility are exact only in the specified game; the reciprocity example demonstrates why that impossibility cannot be generalized to every form of clientelism.
\begin{thebibliography}{9}
\bibitem{fs} F. Finan and L. Schechter (2012).
\emph{Vote-Buying and Reciprocity}, Econometrica 80(2), 863--881.
\url{https://onlinelibrary.wiley.com/doi/10.3982/ECTA9035}.
The primary article motivates reciprocity as a distinct mechanism; the game and thresholds above are explicit restricted constructions.
\end{thebibliography}

\end{document}
